Large language models (LLMs) frequently generate factually incorrect content---hallucinations---despite advances in scale and training. We investigate whether certain hallucinations are \emph{natural}: consistently produced across different model families, robust to question rephrasing, and difficult for models to recognize as errors. Testing 100 questions from \truthfulqa across \gptfour, \gptthree, \claude, and \llama, we find that 7\% of questions cause failures in three or more models simultaneously. These \naturalhall exhibit 57\% robustness to paraphrasing, and models recognize only 24\% of their own errors on these questions when asked directly. We also find that 20\% of errors made by \gptthree persist in \gptfour, suggesting that older model failures partially predict newer model vulnerabilities. Our findings indicate that \naturalhall represent systematic blindspots arising from shared training data biases rather than individual model quirks, with implications for hallucination mitigation strategies: self-critique approaches will fail on these cases, and addressing them requires understanding common biases across training corpora rather than model-specific interventions.
